\section{Distances between dyadic layers}
\label{sec:area_law_layer_distance}

Retain the cells, neighborhoods and layers of
Definition~\ref{def:area_law_dyadic_layers}. The origin $o\in\mathbb R^2$ is
arbitrary, $Z\subset\mathbb Z^2$ is finite, and $C_0\ge0$ is an integer.
All scale indices are nonnegative integers. Distances in the plane are taken
in the sup metric, so $d_\infty(x,y)=\max\{|x_1-y_1|,|x_2-y_2|\}$.

\begin{theorem}[Distances between closed cells]
    \label{thm:area_law_dyadic_cell_distance}
    \lean{TNLean.PEPS.AreaLaw.Geometry.dyadicCell_dist_le,
        TNLean.PEPS.AreaLaw.Geometry.dyadicCell_dist_lower}
    \leanok
    \uses{def:area_law_dyadic_cells}
    Let $u,v\in\mathbb Z^2$, $x\in\overline{C_{k,u}}$, and
    $y\in\overline{C_{k,v}}$. If $|u_i-v_i|\le C_0$ for both coordinates
    $i\in\{1,2\}$, then
    \begin{align}
        d_\infty(x,y) & \le (C_0+1)2^k.
    \end{align}
    If $|u_i-v_i|>C_0$ in at least one coordinate, then
    \begin{align}
        C_0 2^k & \le d_\infty(x,y).
    \end{align}
\end{theorem}
\begin{proof}\leanok
    \uses{thm:area_law_dyadic_closures}
    Each coordinate of a closed cell lies in its closed interval of length
    $2^k$. Indices differing by at most $C_0$ therefore give coordinate
    distances at most $(C_0+1)2^k$. If two integer indices differ by more
    than $C_0$, they differ by at least $C_0+1$; the corresponding closed
    intervals are separated by at least $C_0 2^k$. Taking the maximum of
    the two coordinate distances proves both bounds.
\end{proof}

% Source: geometry:layer-distance in Section 11.
\begin{theorem}[Distance from a layer to the endpoints]
    \label{thm:area_law_dyadic_layer_distance}
    \lean{TNLean.PEPS.AreaLaw.Geometry.dyadicLayer_dist_lower,
        TNLean.PEPS.AreaLaw.Geometry.dyadicLayer_exists_dist_le,
        TNLean.PEPS.AreaLaw.Geometry.dyadicLayer_infDist_bounds}
    \leanok
    \uses{def:area_law_dyadic_layers}
    For every $x\in\overline{D_k}$ and every $z\in Z$,
    \begin{align}
        C_0 2^k & \le d_\infty(x,z).
    \end{align}
    For every $x\in\overline{D_k}$ there is a point $z_*\in Z$ such that
    \begin{align}
        d_\infty(x,z_*) & \le 2(C_0+1)2^k.
    \end{align}
    In particular, the existence of such an $x$ implies $Z\ne\varnothing$.
    Writing $d_\infty(x,Z)=\inf_{z\in Z}d_\infty(x,z)$, one obtains
    \begin{align}
        C_0 2^k & \le d_\infty(x,Z) \le 2(C_0+1)2^k.
    \end{align}
    This is the layer-distance estimate in
    \cite[Section~11]{OpenAI2026AreaLaw}. Empty layers are allowed.
\end{theorem}
\begin{proof}\leanok
    \uses{thm:area_law_dyadic_closures,thm:area_law_dyadic_indices,
        thm:area_law_dyadic_cell_distance}
    The closure of the layer is a finite union of closed scale-$k$ cells.
    A layer index lies outside $P_k$, so it differs from every occupied
    index in at least one coordinate by more than $C_0$. Applying the
    closed-cell lower bound to the cell containing $z$ proves the first
    inequality.

    The inclusion $D_k\subseteq N_{k+1}$ persists after taking closures.
    A closed cell in $N_{k+1}$ has an index within $C_0$ of an occupied
    index. Choose a point $z_*$ of $Z$ in the latter cell and apply the
    closed-cell upper bound at scale $k+1$. Its side length is $2^{k+1}$,
    giving the second inequality. The infimum is bounded below by the
    first inequality and above by the distance to $z_*$.
\end{proof}

% Source: geometry:nonadjacent in Section 11.
\begin{theorem}[Separation of nonadjacent layers]
    \label{thm:area_law_dyadic_nonadjacent}
    \lean{TNLean.PEPS.AreaLaw.Geometry.dyadicLayer_dist_separation,
        TNLean.PEPS.AreaLaw.Geometry.dyadicLayer_dist_nonadjacent}
    \leanok
    \uses{def:area_law_dyadic_layers}
    For every pair of scales $k,h$, every $x\in\overline{D_k}$, and every
    $y\in\overline{D_h}$,
    \begin{align}
        C_0 2^h-2(C_0+1)2^k & \le d_\infty(x,y).
    \end{align}
    If $h\ge k+2$, then
    \begin{align}
        \frac{C_0-1}{2}2^h & \le d_\infty(x,y).
    \end{align}
    This is the nonadjacent-layer estimate in
    \cite[Section~11]{OpenAI2026AreaLaw}. The last lower bound is positive
    under the source's choice $C_0\ge2$. Both inequalities hold for every
    $C_0\ge0$, including when the stated lower bound is nonpositive.
\end{theorem}
\begin{proof}\leanok
    \uses{thm:area_law_dyadic_layer_distance}
    Choose $z\in Z$ with $d_\infty(x,z)\le2(C_0+1)2^k$. The endpoint
    lower bound gives $d_\infty(y,z)\ge C_0 2^h$. The triangle inequality
    proves the first assertion. If $h\ge k+2$, then $4\cdot2^k\le2^h$,
    and hence
    \begin{align}
        C_0 2^h-2(C_0+1)2^k & \ge \frac{C_0-1}{2}2^h.
    \end{align}
    This proves the second assertion.
\end{proof}
